
\documentclass[9pt,letterpaper]{article}
\usepackage[margin=0.42in]{geometry}
\usepackage[T1]{fontenc}
\usepackage{lmodern}
\usepackage{enumitem}
\usepackage{xcolor}
\definecolor{accent}{HTML}{1F4E79}
\usepackage[colorlinks=true,linkcolor=accent,urlcolor=accent,citecolor=accent]{hyperref}
\usepackage{titlesec}
\usepackage{microtype}
\pagestyle{empty}
\setlength{\parindent}{0pt}
\setlength{\parskip}{0pt}
\setlist[itemize]{leftmargin=1.1em,itemsep=0.55pt,topsep=0.8pt,parsep=0pt,partopsep=0pt}
\titleformat{\section}{\normalsize\bfseries\color{accent}}{}{0em}{}[\titlerule]
\titlespacing*{\section}{0pt}{5pt}{1.5pt}
\newcommand{\subgroup}[1]{\vspace{1.8pt}{\small\bfseries\itshape\color{accent} #1}\vspace{0.4pt}}
\newcommand{\entry}[1]{\vspace{1.6pt}{\bfseries #1}\vspace{0.3pt}}
\begin{document}

\begin{center}{\Large\bfseries Ziyu Zhao}\end{center}\vspace{-5pt}
\begin{center}\footnotesize Email: ziyuzhao2027@163.com | Tel: +86-18439471396 | GitHub: https://github.com/12sqawdwq\end{center}\vspace{-4pt}
\section*{Education}
\entry{Xi'an Jiaotong-Liverpool University | Sep 2023 - Present}
BEng in Microelectronics Science and Technology | Average Grade: 82.5\\
\section*{Internship Experience}
\entry{Shenzhen InnoX Academy | Jul 2025 - Present}
Embedded / Algorithm Development Intern\\
\begin{itemize}
\item Developed embedded modules for a medical intraocular-pressure measurement device in the LanZhi Medical team, covering Bluetooth communication and closed-loop motor control.
\item Transitioned into algorithm development for probe-driven ocular biomechanical response measurement; built finite-element models in ANSYS to analyze non-ideal measurement conditions, including eccentric probe advancement, subject-specific biomechanical variance, and operational errors.
\item Developed control, error-compensation, and sampling strategies for probe-offset advancement, biomechanical variance, and measurement-operation disturbances.
\end{itemize}
\entry{ByteDance Volcano Engine Data Department | Dec 2025 - May 2026}
Technical Operations Intern\\
\begin{itemize}
\item Handled technical operations for community-facing AI products, including user Q\&A, issue triage, and feedback consolidation for product and technical teams.
\item Conducted horizontal benchmark comparisons of large-model capabilities across representative tasks; recorded performance differences, limitations, and task-specific operating conditions.
\item Analyzed competing LLM agent products and vibe-coding workflows, covering product positioning, tool-chain structure, developer interaction flow, and feature coverage.
\end{itemize}
\entry{DeepBlue Technology (Shanghai) | Apr 2024 - Jun 2025}
Embedded Software Engineer\\
\begin{itemize}
\item Implemented gait-control algorithms and motor-control hardware interfaces for a wheeled biped robot.
\item Used Raspberry Pi 5 for path planning, SLAM-based mapping, and real-time obstacle avoidance with a depth camera; designed CAN-based communication with embedded controllers.
\item Applied IMU-based sensor fusion for localization and motion-state estimation.
\item Delivered technical presentations on DeepBlue AI and robotics products in technology forums and campus lectures.
\end{itemize}
\entry{Sanyi 3D Technology (Qingdao) | Jan 2024 - Mar 2024}
Embedded Full-Stack Engineer\\
\begin{itemize}
\item Designed and assembled large-scale FDM 3D printers to support product iteration.
\item Used SolidWorks to design and simulate delta-printer motion systems, focusing on stepper-motor actuation and belt-drive transmission.
\item Deployed Marlin firmware for sensor and actuator control; implemented PID-based motor control and sensor integration.
\end{itemize}
\section*{Publication}
\begin{itemize}
\item Z. Zhao, H. S. Gan, Y. Cao, Z. A. Ocal, and I. O. Koska, "PAD-Former: A Two-Stage PI-RADS Classification Framework with Multi-Modal Anomaly Diffusion and Cross-Attention," in \textit{Proc. 2025 IEEE International Conference on Bioinformatics and Biomedicine (BIBM)}, 2025, pp. 6162-6169, doi: \href{https://doi.org/10.1109/BIBM66473.2025.11356522}{10.1109/BIBM66473.2025.11356522}. (BIBM conference paper; CCF-B)
\end{itemize}
\section*{Project Experience}
\subgroup{AI \& Engineering Tools}
\entry{ForgeSpec Studio: AI-assisted Parametric CAD Generator | GitHub: \href{https://github.com/12sqawdwq/ForgeSpec-Studio}{ForgeSpec-Studio} | 2026}
\begin{itemize}
\item Built an AI-assisted engineering tool that turns natural-language requirements into source-first parametric CAD packages with browser preview and STEP / STL / Python / metadata export.
\item Designed a FastAPI backend with a deterministic AssemblySpec pipeline, CadQuery-based CAD generation, provider support for Zhipu / BigModel and Gemini, and fallback templates for standard parts.
\item Added bilingual UI, progress logs, JSON highlighting, one-click downloads, and standards-backed expansion for common fasteners and mechanical components.
\end{itemize}
\entry{Sky Monitor API \& Interactive Star-Map Wiki | GitHub: \href{https://github.com/12sqawdwq/sky_monitor_api}{sky\_monitor\_api} | 2026}
\begin{itemize}
\item Built a FastAPI backend for a location-aware stargazing assistant, returning structured JSON for visible stars, planets, and deep-sky objects based on GPS coordinates and observation time.
\item Integrated HYG star catalog data, Skyfield / JPL DE421 ephemerides, SQLite catalog queries, and Messier / NGC object metadata to support astronomy calculations and target filtering.
\item Developed an interactive Wiki demo with polar star-chart rendering, city lookup, zoom / rotation controls, file upload, and LLM-ready reasoning context for observation recommendations.
\end{itemize}
\subgroup{Embedded Systems \& Robotics}
\entry{ROS2-based Autonomous Vehicle Control System | Fudan University Intelligent Perception \& Autonomous System Lab | Jun 2024 - Aug 2024}
\begin{itemize}
\item Developed a LiDAR-based perception-fusion module for environment sensing and obstacle avoidance on NVIDIA Jetson NX.
\item Implemented SLAM for localization and mapping, A* path planning in ROS2, and PID-based speed and steering control.
\item Designed ROS2 navigation and control pipeline modules integrating perception, planning, and low-level actuation.
\end{itemize}
\entry{OpenMV-based High-Speed Target Detection System | GitHub: \href{https://github.com/12sqawdwq/openmv_green_tracking}{openmv\_green\_tracking} | Jun 2025 - Aug 2025}
\begin{itemize}
\item Implemented a real-time multi-color target detection system using OpenMV and STM32H7, achieving 120-200 FPS in dynamic conditions.
\item Developed vision-based distance estimation through camera calibration and designed a high-speed SPI communication interface for embedded data transfer.
\end{itemize}
\entry{STM32 20x4 LCD Graphics Engine | GitHub: \href{https://github.com/12sqawdwq/mes204_plus}{mes204\_plus} | 2025}
\begin{itemize}
\item Built a graphics engine on STM32F446xx for an HD44780-compatible 20x4 character LCD, using an optimized 8-bit parallel LCD driver and CMake / ARM GCC build workflow.
\item Implemented graphics primitives, mathematical visual effects, a wireframe 3D renderer, particle effects, and character-video playback under highly constrained display resources.
\item Structured the repository with bilingual documentation, firmware architecture notes, build instructions, and visual demonstration materials.
\end{itemize}
\entry{STM32 LCD1602 Real-Time Game System | GitHub: \href{https://github.com/12sqawdwq/MES204_lcd}{MES204\_lcd} | 2025}
\begin{itemize}
\item Built an embedded real-time T-Rex runner game on STM32F446xx using C, LCD1602 character display rendering, timer interrupts, and EXTI-based low-latency button input.
\item Designed a constrained rendering pipeline using LCD1602 CGRAM resources and a non-blocking state machine to support responsive animation under limited display memory.
\item Packaged the project with hardware documentation, firmware architecture notes, and reproducible STM32 development setup.
\end{itemize}
\subgroup{Digital Design \& Chip-Level Systems}
\entry{AXI-Stream Header Insertion RTL Module | GitHub: \href{https://github.com/12sqawdwq/verilog_code_test}{verilog\_code\_test} | 2025}
\begin{itemize}
\item Designed a Verilog AXI-Stream module for packet-header insertion with valid / ready handshaking, keep / last byte-valid signals, and configurable data width.
\item Implemented byte-count tracking, data buffering, and insertion-state control to align inserted header data with streaming payload transfer.
\item Documented module parameters, ports, internal registers, and output-transfer logic for digital-design review and reuse.
\end{itemize}
\subgroup{Computer Systems \& Developer Tools}
\entry{STM32 CLI Toolkit for Remote Embedded Development | GitHub: \href{https://github.com/12sqawdwq/stm32-cli-toolkit}{stm32-cli-toolkit} | 2025}
\begin{itemize}
\item Developed a lightweight Bash / Zsh toolkit wrapping OpenOCD, CMake, and ARM GCC workflows for STM32 development on Linux and headless environments.
\item Implemented commands for automatic build-artifact detection, multi-series flashing, parallel compilation, and USB reset recovery in STM32 bring-up workflows.
\end{itemize}
\subgroup{Mathematical Modeling \& Biomedical Systems}
\entry{CAE Simulation Validation Platform | GitHub: \href{https://github.com/12sqawdwq/ANSYS_skill}{ANSYS\_skill} | 2026}
\begin{itemize}
\item Built a simulation-validation workflow around Discovery / SpaceClaim geometry handoff, Workbench project setup, Mechanical meshing, load definition, solving, and result evaluation.
\item Defined FEM/CFD validation checks covering solved-state evidence, contour and legend inspection, and non-empty Min / Max result verification.
\end{itemize}
\entry{STM32-based Smart Pressure-Sensing Cushion System | East China Normal University Intelligent Ultrasound Research Center | Feb 2024 - May 2024}
\begin{itemize}
\item Designed a high-resolution pressure-sensor array system for real-time sitting-posture monitoring.
\item Used KNN regression to identify abnormal postures, including slouching and pelvic tilt, and estimate spinal pathological features.
\item Visualized real-time pressure distributions and posture analytics in MATLAB through heatmaps and monitoring dashboards.
\end{itemize}
\entry{Multi-channel Electrical Stimulation Control System | East China Normal University Intelligent Ultrasound Research Center | Nov 2023 - Feb 2024}
\begin{itemize}
\item Developed a cross-platform Flutter mobile control system and aligned communication between the hardware layer and back-end protocol processing.
\item Designed system architecture, including front-end UI/UX, communication protocol handling, device-control logic, and integration of I2C, CAN, and Bluetooth modules.
\item Developed embedded firmware for STM32F103/F407 and ESP32 prototypes for synchronized multi-channel stimulation output.
\item Implemented closed-loop stimulation control based on bio-impedance feedback.
\end{itemize}
\entry{Wearable ECG \& Oxygen Saturation (SpO2) Monitoring Device | Shanghai Key Lab of Multi-Dimensional Information Processing | Nov 2023 - Feb 2024}
\begin{itemize}
\item Created a portable biomedical sensing device using MAX30102 / GH3220 for ECG and SpO2 monitoring.
\item Filtered biomedical signals, extracted physiological parameters, and developed Bluetooth / serial data transmission.
\item Built a user-facing real-time monitoring interface for physiological-data visualization.
\end{itemize}
\section*{Campus Activities}
\begin{itemize}
\item President, XJTLU Science Innovation Association.
\item Campus Ambassador: JetBrains, MATLAB, and ZhipuAI.
\end{itemize}
\section*{Skills}
\begin{itemize}
\item Programming \& HDL: C/C++, Python, MATLAB, Verilog, JavaScript, HTML/CSS
\item Embedded \& Robotics: STM32, ESP32, Raspberry Pi, OpenMV, ROS2, CAN, I2C, SPI, Bluetooth, PID control, sensor fusion, SLAM
\item AI / Vision / Agents: PyTorch, TensorFlow, OpenCV, large-model benchmarking, LLM agent analysis, prompt and workflow evaluation
\item AI Coding Tools: Claude Code, Codex, Cursor, GitHub Copilot, Windsurf
\item Engineering Tools: ANSYS, SolidWorks, CadQuery, GNU Toolchain, OpenOCD, CMake, Git
\end{itemize}
\end{document}
